\begin{thebibliography}{99}
\setlength{\itemsep}{0pt}
\bibitem{unet} Ronneberger O, Fischer P, Brox T. U-Net: Convolutional Networks for Biomedical Image Segmentation[C]. MICCAI, 2015: 234-241.
\bibitem{unet3d} Cicek O, Abdulkadir A, Lienkamp S S, et al. 3D U-Net: Learning Dense Volumetric Segmentation from Sparse Annotation[C]. MICCAI, 2016: 424-432.
\bibitem{attention_unet} Oktay O, Schlemper J, Le Folgoc L, et al. Attention U-Net: Learning Where to Look for the Pancreas[J]. arXiv preprint arXiv:1804.03999, 2018.
\bibitem{nnunet} Isensee F, Jaeger P F, Kohl S A A, et al. nnU-Net: A Self-configuring Method for Deep Learning-based Biomedical Image Segmentation[J]. Nature Methods, 2021, 18(2): 203-211.
\bibitem{kamnitsas2017} Kamnitsas K, Ledig C, Newcombe V F J, et al. Efficient Multi-scale 3D CNN with Fully Connected CRF for Accurate Brain Lesion Segmentation[J]. Medical Image Analysis, 2017, 36: 61-78.
\bibitem{chen2021} Chen M, Geng C, Wang D, et al. Deep Learning-based Segmentation of Cerebral Aneurysms in 3D TOF-MRA using Coarse-to-Fine Framework[J]. arXiv preprint arXiv:2110.13432, 2021.
\bibitem{ou2022} Ou C, Qian Y, Chong W, et al. A Deep Learning-based Automatic System for Intracranial Aneurysms Diagnosis on Three-dimensional DSA Images[J]. Medical Physics, 2022, 49(11): 7038-7053.
\bibitem{strokenet} Irfan M, Malik K M, Ahmad J, et al. StrokeNet: An Automated Approach for Segmentation and Rupture Risk Prediction of Intracranial Aneurysm[J]. Computerized Medical Imaging and Graphics, 2023, 108: 102271.
\bibitem{xiong2025} Xiong Y, Ye M, Wang P, et al. Rupture Risk Prediction of Intracranial Aneurysm by Using Gene Expression Data Mining and Intelligent Optimization Algorithm[J]. International Journal of Bio-Inspired Computation, 2025, 26(1): 24-34.
\end{thebibliography}
